% ============ 29 非常值层的 Lefschetz 不动点公式 (原书 pp.163-175) ============
\chapter{非常值层的 Lefschetz 不动点公式}\label{sec:29}
\index{Lefschetz 不动点公式}
\index{$\zeta$ 函数}
\index{Artin L 级数}
\index{完美复形}
\index{非交换迹}
\index{反变表示}

我们需要非完备簇，非常值层与有限环上模层的 Lefschetz 不动点公式.这些推广中的每一个都引起问题，我们现在讨论它们.

\textbf{非完备簇} 设 $U$ 是代数闭域 $k$ 上完备非奇异簇 $X$ 的开子集，并假定 $U$ 在 $X$ 中的补集 $Z$ 非奇异.设 $\Lambda = \mathbb{Q}_{\ell}$.由正合序列 (\ref{prop:8.15})
\[
0 \to j_{!}\Lambda \to \Lambda \to i_{*}\Lambda \to 0
\]
得到长正合序列
\[
\cdots \to H^{r}_{c}(U, \Lambda) \to H^{r}(X, \Lambda) \to H^{r}(Z, \Lambda) \to \cdots.
\]
设 $\varphi \colon X \to X$ 是有限正则映射.若 $\varphi$ 保持分解 $X = U \cup Z$，则它作用在该复形上，从而
\[
\sum_{r} (-1)^{r} \operatorname{Tr}(\varphi|H^{r}(X)) = \sum_{r} (-1)^{r} \operatorname{Tr}(\varphi|H^{r}_{c}(U)) + \sum_{r} (-1)^{r} \operatorname{Tr}(\varphi|H^{r}(Z)).
\]
设 $\Gamma^{\varphi}$ 表示被 $\varphi$ 固定的 $\Gamma$ 的闭点集.若
\[
\#X^{\varphi} = \#U^{\varphi} + \#Z^{\varphi},
\]
则作为 $X$ 与 $Z$ 的不动点公式之差，得到不动点公式
\[
\#U^{\varphi} = \sum (-1)^{r} \operatorname{Tr}(\varphi|H^{r}_{c}(U, \Lambda)).
\]
然而，如下一个例子所示，这个论证会误导.

\begin{example}\label{ex:29.1}
考虑 (\ref{ex:25.2}) 的映射
\[
\varphi \colon \mathbb{P}^{1} \to \mathbb{P}^{1}, \qquad (x_0 \colon x_1) \mapsto (x_0 + x_1 \colon x_1).
\]
设 $U = \{(x_0 \colon x_1) \mid x_1 \neq 0\} \cong \mathbb{A}^{1}$，$Z = \{(1 \colon 0)\} = \{1\}$.则对 $r \neq 2$，$H^{r}_{c}(U, \mathbb{Q}_{\ell}) = 0$，且 $\varphi$ 在 $H^{2}_{c}(U, \mathbb{Q}_{\ell}) \cong H^{2}(\mathbb{P}^{1}, \mathbb{Q}_{\ell})$ 上按恒等映射作用.因此，尽管 $\varphi|_{U}$ 没有不动点($\varphi$ 在 $U$ 上按 $x \mapsto x + 1$ 作用), $\sum (-1)^{r} \operatorname{Tr}(\varphi|H^{r}_{c}(U, \mathbb{Q}_{\ell})) = 1$---Lefschetz 不动点公式对 $U$ 与 $\varphi|_{U}$ 不成立.

问题是 $1$ 作为 $\varphi|_{Z}$ 的不动点重数为 1，但作为 $\varphi$ 的不动点重数为 2.因此，按重数计数时，等式
\[
\#X^{\varphi} \neq \#U^{\varphi} + \#Z^{\varphi}.
\]
于是，只有当映射延拓为簇的完备化上的映射，且它在簇的补集上只有简单不动点时，才应期望非完备簇有 Lefschetz 不动点公式---事实上，上面的论证确实正确地表明在这种情况下(常数系数)得到不动点公式.
\end{example}

\textbf{非常值层} 设 $X$ 是代数闭域 $k$ 上的完备簇，$E$ 是 $X$ 上 $\Lambda$-模的局部常值层.正则映射 $\varphi \colon X \to X$ 定义映射 $H^{r}(X, E) \to H^{r}(X, \varphi^{*}E)$---这是不同向量空间之间的映射，因此它的迹未定义.为了有迹，还需要同态 $\varphi_E \colon \varphi^{*}E \to E$.这样的对 $(\varphi, \varphi_E)$ 定义映射
\[
H^{r}(X, E) \to H^{r}(X, \varphi^{*}E) \to H^{r}(X, E),
\]
我们把它的复合记为 $(\varphi, \varphi_E)^{*}$.对 $X$ 的每个闭点 $x$，$\varphi_E$ 在茎上定义映射：
\[
\varphi_x \colon (\varphi^{*}E)_x \to E_x
\]
\[
\parallel
\]
\[
E_{\varphi(x)}.
\]
当 $x$ 是 $\varphi$ 的不动点时，这变成
\[
\varphi_x \colon E_x \to E_x.
\]
我们可能希望在适当假设下，有公式
\[
\sum_{x \in X^{\varphi}} \operatorname{Tr}(\varphi_x|E_x) = \sum (-1)^{r} \operatorname{Tr}((\varphi, \varphi')^{*}|H^{r}(X, E)),
\]
其中 $X^{\varphi}$ 是 $\varphi$ 的闭不动点集.

当 $E$ 常值，设由群 $E$ 定义时，$\varphi^{*}E$ 也是由 $E$ 定义的常值层.因此，在这种情况下映射 $\varphi_E$ 是恒等映射.例如，若 $E = \mathbb{Q}_{\ell}$，公式变成
\[
\#X^{\varphi} = \sum (-1)^{r} \operatorname{Tr}(\varphi|H^{r}(X, \mathbb{Q}_{\ell})),
\]
其中 $\#X^{\varphi}$ 是不计重数的不动点数.

\begin{remark}\label{rem:29.2}
设 $X$ 是有限集，把 $X$ 看作离散拓扑空间.给出 $X$ 上有限维 $\mathbb{Q}$-向量空间的层 $E$ 就是给出由 $X$ 的元素指标化的有限维向量空间族 $(E_x)_{x \in X}$.如上的映射对 $(\varphi, \varphi_E)$ 是集合的映射 $\varphi \colon X \to X$ 与由 $X$ 的元素指标化的映射族 $\varphi_x \colon E_{\varphi(x)} \to E_x$.映射 $(\varphi, \varphi_E)^{*} \colon H^{0}(X, E) \to H^{0}(X, E)$ 是映射 $\varphi_x \colon E_{\varphi(x)} \to E_x$ 的直和.显然，除非 $\varphi(x) = x$(此时贡献 $\operatorname{Tr}(\varphi_x)$)，$\varphi_x$ 对 $\operatorname{Tr}((\varphi, \varphi_E)^{*}|H^{0}(X, E))$ 没有贡献.因此公式在这种情况下成立.
\end{remark}

\textbf{有限系数环} 若 $\Lambda$ 不是域，则 $H^{r}(X, E)$ 可能不是自由 $\Lambda$-模，这时自同态的迹未定义.由于这个问题直到定理~\ref{thm:29.4} 的证明中才出现，我们推迟讨论它.

\subsection*{定理的陈述}

\begin{lemma}\label{lem:29.3}
设 $X_0$ 是 $\mathbb{F}_q$ 上的簇，$E_0$ 是 $X_{0\mathrm{et}}$ 上的层.设 $E$ 是 $E_0$ 在 $X$ 上的逆像.则有 $X$ 上层 的典范同态 $F_E \colon F^{*}E \to E$.
\end{lemma}

\begin{proof}
(梗概) 我们只在 $E_0$ 是 $X_0$ 上的簇 $\varphi'_0 \colon E_0 \to X_0$ 定义的层时描述 $F_E$.则 $E$ 是 $\varphi \colon E \to X$ 定义的层.考虑图
\[
\begin{tikzcd}
E \arrow[d, "\varphi"'] \arrow[rr, "F"] & & E \arrow[d, "\varphi"] \\
X \arrow[rr, "F"'] & & X
\end{tikzcd}
\]
给出 $\Gamma(X, E)$ 的元素就是给出满足 $\varphi \circ s = \mathrm{id}$ 的截面 $s \colon X \to E$；给出 $\Gamma(X, F^{*}E)$ 的元素就是给出满足 $\varphi \circ s' = F$ 的映射 $s' \colon X \to E$.映射 $\Gamma(X, E) \to \Gamma(X, F^{*}E)$ 把 $s$ 送到 $s' \overset{\mathrm{def}}{=} s \circ F$，而 $F_E \colon \Gamma(X, F^{*}E) \to \Gamma(X, E)$ 把 $s'$ 送到满足 $F \circ s'' = s'$ 的唯一 $s''$.由于所有可构造层都可以由代数空间表示，该论证适用于所有这样的层.
\end{proof}

$\mathbb{F}_q$ 上的非奇异曲线 $U_0$ 可以(本质唯一地)嵌入 $\mathbb{F}_q$ 上的完备非奇异曲线 $X_0$.Frobenius 映射 $F \colon X \to X$ 保持 $U$，它在 $X \setminus U$ 中的不动点重数为 1.因此，如下定理至少是可信的.

\begin{theorem}\label{thm:29.4}
设 $U_0$ 是 $\mathbb{F}_q$ 上的非奇异曲线，$E_0$ 是 $U_0$ 上 $\mathbb{Q}_{\ell}$-向量空间的局部常值层.则
\[
\sum_{x \in U^{F}} \operatorname{Tr}(F_x|E_x) = \sum (-1)^{r} \operatorname{Tr}(F|H^{r}_{c}(U, E)).
\]
左边的和遍历被 $F$ 固定的 $U$ 的闭点(它们与 $U_0(\mathbb{F}_q)$ 的元素自然一一对应)，$E_x$ 是 $E$ 在 $x$(看作几何点)处的茎，$F_x$ 是由 $F_E$ 在茎上诱导的映射.右边，$F$ 是由 $(F, F_E)$ 诱导的映射.
\end{theorem}

\begin{example}\label{ex:29.5}
设 $\varphi \colon Y_0 \to U_0$ 是 $U_0$ 上完备非奇异曲线的族，$E_0 = R^{1}\varphi_{*}\mathbb{Q}_{\ell}$.对 $x \in U_0(\mathbb{F}_q)$，设 $Y_x = \varphi^{-1}(x)$.它是 $\mathbb{F}_q$ 上的曲线，且
\[
\operatorname{Tr}(F_x|E_x) = \operatorname{Tr}(F|H^{1}(Y_x, \mathbb{Q}_{\ell})).
\]
\end{example}

在解释定理~\ref{thm:29.4} 的证明之前，我们讨论一些应用，并用 $\pi_1(U_0)$-模的语言重新解释定理.

\subsection*{\texorpdfstring{局部常值层的 $\zeta$ 函数}{局部常值层的 zeta 函数}}

设 $U_0$ 是 $\mathbb{F}_q$ 上的非奇异曲线，$E_0$ 是 $U_0$ 上 $\mathbb{Q}_{\ell}$-向量空间的局部常值层.定义 $Z(U_0, E_0, t)$ 为
\[
\log Z(U_0, E_0, t) = \sum_{m > 0} \left( \sum_{x \in U^{F^{m}}} \operatorname{Tr}(F^{m}_x|E_x) \right) \frac{t^{m}}{m}.
\]
例如，
\[
Z(U_0, \mathbb{Q}_{\ell}, t) = \exp\left( \sum_{m > 0} N_m \frac{t^{m}}{m} \right) = Z(U_0, t).
\]
与 \S 27 相同的论证(见引理~\ref{lem:27.9}.1)表明
\[
Z(U_0, E_0, t) = \prod_{x \in U_0} \frac{1}{\det(1 - F_x t^{\deg x}|E_x)}.
\]

\begin{theorem}\label{thm:29.6}
用上述记号，
\[
Z(U_0, E_0, t) = \frac{\det(1 - Ft|H^{1}_{c}(U, E))}{\det(1 - Ft|H^{0}_{c}(U, E)) \cdot \det(1 - Ft|H^{2}_{c}(U, E))}.
\]
\end{theorem}

从 29.4 推出该定理与从定理~\ref{thm:25.1} 推出定理~\ref{thm:27.6} 相同.事实上，定理~\ref{thm:29.6} 常称为不动点公式的乘法形式.

\begin{remark}\label{rem:29.7}
当 $U_0$ 仿射时，$H^{0}_{c}(U, E) = 0$($U$ 上局部常值层的截面没有完备子簇上的支集)，因此等式变成
\[
Z(U_0, E_0, t) = \frac{\det(1 - Ft|H^{1}_{c}(U, E))}{\det(1 - Ft|H^{2}_{c}(U, E))}.
\]
\end{remark}

\subsection*{\texorpdfstring{任意簇的 $\zeta$ 函数}{任意簇的 zeta 函数}}

我们像好\footnote{Shimura 簇的研究提示，事实上，对非完备簇这不是正确的定义.}的情形一样定义 $\mathbb{F}_q$ 上任意(可能奇异且不完备的)簇 $X_0$ 的 $\zeta$ 函数 $Z(X_0, t)$：
\[
Z(X_0, t) = \exp\left( \sum_{m \geq 1} N_m \frac{t^{m}}{m} \right) = \prod_{x \in X_0} \frac{1}{1 - t^{\deg x}},
\]
其中 $N_m$ 是 $X_0$ 上坐标在 $\mathbb{F}_{q^m}$ 中的点的个数，乘积遍历 $X_0$ 的闭点.

\begin{theorem}\label{thm:29.8}
设 $X_0$ 是 $\mathbb{F}_q$ 上的簇.则
\[
Z(X_0, t) = \frac{P_1(X_0, t) \cdots P_{2d-1}(X_0, t)}{P_0(X_0, t) \cdots P_{2d}(X_0, t)},
\]
其中
\[
P_r(X, t) = \det(1 - Ft|H^{r}_{c}(X, \mathbb{Q}_{\ell})).
\]
\end{theorem}

关于定理~\ref{thm:29.8} 由定理~\ref{thm:29.4} 推出的证明，见 EC pp.289--298.

定理表明 $Z(X_0, t) \in \mathbb{Q}(t)$，但目前还不知道 $P_r(X, t)$ 是否与 $\ell$ 无关，甚至不知道它们是否有 $\mathbb{Q}$ 中的系数.

\subsection*{\texorpdfstring{用 $\pi_1(U)$-模重新解释；Frobenius 元素}{用 pi_1(U)-模重新解释；Frobenius 元素}}

现在我们假定 $U_0$ 是 $\mathbb{F}_q$ 上的仿射非奇异曲线，并设 $K$ 是 $U$ 的函数域 $\mathbb{F}_q(U_0)$.固定 $K$ 的代数闭包 $\Omega$.则 $U_0 = \operatorname{Specm} A$，其中 $A = \Gamma(U_0, \mathcal{O}_{U_0}) \subset K$ 是 Dedekind 整环，我们就把 $\pi_1(U_0)$ 等同于 $\operatorname{Gal}(K^{\mathrm{un}}/K)$，其中 $K^{\mathrm{un}}$ 是 $A$ 的所有素理想上都在 $K$ 上非分歧的 $\Omega$ 的子域的并.注意 $K \cdot \mathbb{F} \subset K^{\mathrm{un}}$；我们把 $\pi_1(U)$ 等同于 $\operatorname{Gal}(K^{\mathrm{un}}/K \cdot \mathbb{F})$.

设 $\Lambda = \mathbb{Z}/\ell^{n}\mathbb{Z}$ 或 $\mathbb{Q}_{\ell}$.

\begin{proposition}\label{prop:29.9}
设 $E$ 是代数闭域 $k$ 上仿射曲线 $U$ 上 $\Lambda$-模的局部常值层，$E$ 是对应的 $\pi_1(U)$-模.则
\[
H^{0}(U, E) \cong E^{\pi_1(U)} \qquad H^{0}_{c}(U, E) = 0
\]
\[
H^{1}(U, E) \cong H^{1}(\pi_1(U), E) \qquad H^{1}_{c}(U, E) \cong H^{1}(\pi_1(U), \mathcal{L}E(1))^{\vee}
\]
\[
H^{2}(U, E) = 0 \qquad H^{2}_{c}(U, E) \cong E_{\pi_1(U)}(-1).
\]
这里 $E^{\pi_1(U)}$ 是 $\pi_1(U)$ 平凡作用的最大的 $E$ 的子模，$E_{\pi_1(U)}$ 是 $\pi_1(U)$ 平凡作用的最大的 $E$ 的商模.
\end{proposition}

\begin{proof}
关于 $H^{r}(U, E)$ 的论断见 \S 14.关于 $H^{r}_{c}(U, E)$ 的论断由对偶推出.例如
\[
H^{2}_{c}(U, E) = H^{0}(U, \mathcal{L}E(1))^{\vee} = (\mathcal{L}E(1)^{\pi_1})^{\vee} = ((\mathcal{L}E(1))^{\vee})_{\pi_1} = (E(-1))_{\pi_1} = E_{\pi_1}(-1).
\]
\end{proof}

域的同构 $\alpha \colon k \to k'$ 定义 $k$-簇与 $k'$-簇之间的一一对应 $X \leftrightarrow X'$，平展映射与平展覆盖在它下对应.因此它也定义 $X_{\mathrm{et}}$ 上层与 $X'_{\mathrm{et}}$ 上层之间的一一对应 $E \leftrightarrow E'$，且 $H^{r}(X, E) \cong H^{r}(X', E')$.

我们把这一评注应用于 $\varphi \colon \mathbb{F} \to \mathbb{F}$(Frobenius 自同构 $x \mapsto x^{q}$).若 $X$ 与 $E$ 来自 $\mathbb{F}_q$ 上的对象 $X_0$ 与 $E_0$，则 $X' = X$ 且 $E' = E$.因此，对 $\mathbb{F}_q$ 上的任何簇 $X_0$ 与 $X_0$ 上的层 $E_0$，$\varphi$ 定义 $H^{r}(X, E)$ 的自同构(对紧支集上同调类似).

\begin{proposition}\label{prop:29.10}
对 $\mathbb{F}_q$ 上的任何簇 $X_0$ 与 $X_0$ 上的层 $E_0$，由 $(F, F_E)$ 定义的 $H^{r}(X, E)$ 的自同构与 $\varphi$ 定义的互逆(对紧支集上同调类似).
\end{proposition}

证明省略，但我建议读者对 $\mathbb{A}^{1}$ 上的层 $\mathbb{G}_{a}$ 验证该命题.为此，必须验证映射
\[
H^{0}(\mathbb{A}^{1}, \mathbb{G}_{a}) \to H^{0}(\mathbb{A}^{1}, F^{*}\mathbb{G}_{a}) \xrightarrow{F_{\mathbb{G}_{a}}} H^{0}(\mathbb{A}^{1}, \mathbb{G}_{a})
\]
把
\[
P(T) \mapsto P(T^{q}) \mapsto P(T^{q})^{\frac{1}{q}} = (\varphi^{-1}P)(T).
\]
这里 $P(T) \in \mathbb{F}[T]$ 看作 $\mathbb{A}^{1}$ 上正则函数的环，$\varphi^{-1}$ 作用在 $P$ 的系数上.

\begin{remark}\label{rem:29.11}
回忆如下约定.

当群 $G$ 左作用于向量空间 $V$ 时，它按规则
\[
(gf)(v) = f(g^{-1}v), \qquad f \in \check{V}, \ g \in G, \ v \in V
\]
左作用于 $V$ 的对偶 $\check{V}$.这是定义 $G$ 在 $\check{V}$ 上的左作用的唯一自然方式.$G$ 在 $\check{V}$ 上的这个表示称为原表示的反变表示(contragredient representation，亦称对偶表示).

当 $k$-代数 $R$ 左作用于 $k$-向量空间 $V$ 时，$R$ 按规则
\[
(fr)(v) = f(rv), \qquad f \in \check{V}, \ r \in R, \ v \in V
\]
作用于 $\check{V}$.$R$ 在 $V$ 上的左作用变成 $\check{V}$ 上的右作用---一般地，没有定义 $R$ 在 $\check{V}$ 上的左作用的自然方式.

注意，若 $\gamma \in G$ 在 $V$ 上有特征值 $a, b, \ldots$，则它在 $\check{V}$ 上有特征值 $a^{-1}, b^{-1}, \ldots$.然而，若 $\gamma \in R$ 在 $V$ 上有特征值 $a, b, \ldots$，则它在 $\check{V}$ 上有特征值 $a, b, \ldots$.当 $\gamma$ 既可以看作环的元素又可以看作群的元素时，这可能导致混淆.

考虑 $\mathbb{F}_q$ 上的椭圆曲线 $E_0$.数论学家把``Frobenius 的特征值''定义为 $\varphi \in \operatorname{Gal}(\mathbb{F}/\mathbb{F}_q)$ 在 $T_{\ell}E$ 上的特征值.但 $T_{\ell}E$ 是 $H^{1}(E, \mathbb{Z}_{\ell})$ 的对偶，因此它们也是 $\varphi^{-1}$ 在 $H^{1}(E, \mathbb{Z}_{\ell})$ 上的特征值.按命题~\ref{prop:29.10}，它们也是 $F$ 在 $H^{1}(E, \mathbb{Z}_{\ell})$ 上的特征值.于是，幸运的是，数论学家与几何学家一致.

$E$ 的自同态环也含 Frobenius 元素 $\pi$.它如何作用？按定义，$\pi = F$，因此它在 $H^{1}(E, \mathbb{Z}_{\ell})$ 上按 $F$(从而以绝对值 $q^{\frac{1}{2}}$ 的特征值)作用.它也在 $T_{\ell}E$ 上以相同的特征值作用.因为 $H^{1}$ 是反变函子，$\operatorname{End}(E)$ 右作用于它；因为 $T_{\ell}$ 是共变函子，它左作用于 $T_{\ell}$.事实上，$\pi \in \operatorname{End}(E)$ 将是系数为整数，根有绝对值 $q^{\frac{1}{2}}$ 的多项式的根.因此，当 $\operatorname{End}(E)$ 作用于一个空间时，$\pi$ 总是有绝对值 $q^{\frac{1}{2}}$ 的特征值.
\end{remark}

\begin{proposition}\label{prop:29.12}
设 $E_0$ 是对应于 $\pi_1(U_0)$-模 $E$ 的 $U_0$ 上的局部常值层.对任何 $x \in U^{F}$，有同构 $(E_x, F_x) \cong (E, \varphi^{-1}_x)$.这里 $\varphi_x$ 是对应于 $x$ 的 $\operatorname{Gal}(K^{\mathrm{un}}/K)$ 中的 Frobenius 元素($\varphi_x$ 良定义到共轭).
\end{proposition}

命题中的同构非典范，必然如此，因为 $\varphi_x$ 只定义到共轭，但它的存在蕴含
\[
\operatorname{Tr}(F_x|E_x) = \operatorname{Tr}(\varphi^{-1}_x|E).
\]
于是，定理可以改写为：

\begin{theorem}\label{thm:29.13}
设 $E$ 是带 $\pi_1(U_0)$ 连续作用的有限维 $\mathbb{Q}_{\ell}$-向量空间.则
\[
\sum_{x \in U^{F}} \operatorname{Tr}(\varphi^{-1}_x|E) = -\operatorname{Tr}(\varphi^{-1}|H^{1}(\pi_1(U), \mathcal{L}E(1))^{\vee}) + \operatorname{Tr}(\varphi^{-1}|E_{\pi_1(U)}(-1)).
\]
\end{theorem}

\begin{example}\label{ex:29.14}
设 $K$ 是 $\mathbb{F}_q$ 上的单变量函数域，$L$ 是 $K$ 的带 Galois 群 $G$ 的有限 Galois 扩张.设 $\rho \colon G \to \mathrm{GL}(V)$ 是 $G$ 在 $\mathbb{Q}$(或 $\mathbb{Q}$ 的有限扩张)上的有限维向量空间上的表示.设 $A$ 是 $K$ 中分式域为 $K$，在 $L$ 中的整闭包在 $A$ 上非分歧的 Dedekind 整环.定义
\[
L_{\mathrm{Artin}}(s, \rho) = \prod_{\mathfrak{p} \subset A} \frac{1}{\det(1 - \rho(\varphi_x)N\mathfrak{p}^{-s}|V)}, \qquad s \in \mathbb{C}.
\]
这里乘积遍历 $A$ 的非零素理想，$N\mathfrak{p} = (A \colon \mathfrak{p})$.这是 $\rho$ 的 Artin L 级数，只是我省略了一些因子.

设 $U_0 = \operatorname{Specm} A$.则 $G$ 是 $\pi_1(U_0)$ 的商，我们设 $E_0$ 是对应于 $\pi_1(U_0)$-模 $\check{V} \otimes_{\mathbb{Q}} \mathbb{Q}_{\ell}$ 的 $U_0$ 上 $\mathbb{Q}_{\ell}$-向量空间的层.则
\[
L_{\mathrm{Artin}}(s, \rho) = L(U_0, E, q^{-s}).
\]
由定理~\ref{thm:29.13} 的乘法形式推出 $L_{\mathrm{Artin}}(s, \rho) = \frac{P(q^{-s})}{Q(q^{-s})}$，其中 $P(t), Q(t) \in \mathbb{Q}[t]$.此外，$Q(t) = \det(1 - \varphi t|\check{V}^{N_{G}}(-1))$，其中 $N_{G}$ 是 $G$ 在 $\pi_1(U)$ 中的像.因此，$L_{\mathrm{Artin}}(s, \rho)$ 是复变量的亚纯函数，当 $V^{N_{G}} = 0$ 时甚至全纯.
\end{example}

\subsection*{对有限层重述定理}

粗略地说，定理~\ref{thm:29.4} 的证明思路是过渡到 $E$ 在其上常值的有限平展覆盖 $V_0 \to U_0$，然后在 $V_0$ 上应用通常的 Lefschetz 不动点公式.问题是：一般地，$\mathbb{Q}_{\ell}$-向量空间的局部常值层不会在任何有限覆盖上变成常值($\pi_1(U_0)$ 在层茎上的作用不必穿过有限商).因此需要用 $\Lambda$-模层工作，其中 $\Lambda = \mathbb{Z}/\ell^{n}\mathbb{Z}$.但这时 $H^{r}_{c}(U, E)$ 不必是自由 $\Lambda$-模，因此迹的定义不再清楚.尽管如此，我们陈述定理：

\begin{theorem}\label{thm:29.15}
设 $U_0$ 是 $\mathbb{F}_q$ 上的非奇异仿射几何连通曲线，$E_0$ 是 $U_0$ 上 $\Lambda$-模的平坦可构造局部常值层.则
\[
\sum_{x \in U^{F}} \operatorname{Tr}(F_x|E_x) = \sum_{r} (-1)^{r} \operatorname{Tr}(F|H^{r}_{c}(U, E)).
\]
\end{theorem}

对 $E$ 的条件意味着它对应于作为 $\Lambda$-模自由且有限生成的 $\pi_1(U_0)$-模 $E$.我们的首要任务将是解释右边意味着什么.回忆 $U_0$ 上 $\mathbb{Q}_{\ell}$-向量空间的局部常值层 $E$ 是族 $(E_n, f_n)$，其中每个 $E_n$ 是 $\mathbb{Z}/\ell^{n}\mathbb{Z}$-模的平坦可构造局部常值层.定理~\ref{thm:29.4} 通过对每个 $E_n$ 应用定理~\ref{thm:29.15} 并取逆向极限证明(当然，必须检验这是可能的).

\begin{remark}\label{rem:29.16}
注意，我们只需在 $U^{F}$ 为空的情形证明定理~\ref{thm:29.15}.设 $V_0$ 是省略 $U_0$ 的所有 $\mathbb{F}_q$-有理点的 $U_0$ 的开子集.则对 $(V_0, E)$ 与 $(U_0 \setminus V_0, E)$ 的定理~\ref{thm:29.15} 蕴含对 $(U_0, E)$ 的定理(比较本节开头的讨论).因为 $U \setminus V$ 有限，对 $(U_0 \setminus V_0, E)$ 的 Lefschetz 不动点公式本质上是 (\ref{rem:29.2}).
\end{remark}

\subsection*{完美复形}

设 $R$ 是环(如常，Noether).$R$-模的复形
\[
\cdots \to M^{r} \to M^{r+1} \to \cdots
\]
称为完美的，如果它有界(只有有限多个 $M^{r}$ 非零)且每个 $M^{r}$ 是有限生成投射 $R$-模.回忆对交换 Noether 局部环上的模，
\[
\text{投射且有限生成} = \text{自由且有限生成}.
\]
因为这样的模有有限基，可以定义自同态的迹.

\begin{proposition}\label{prop:29.17}
设 $R$ 是交换局部 Noether 环.设 $M^{\bullet}$ 是 $R$-模的复形，$\alpha \colon M^{\bullet} \to M^{\bullet}$ 是 $M^{\bullet}$ 的自同态.对带 $P^{\bullet}$ 完美的任何拟同构 $\psi \colon P^{\bullet} \to M^{\bullet}$，存在 $\beta$ 的自同态 $P^{\bullet}$，唯一到同伦，使得 $\psi \circ \beta = \alpha \circ \psi$：
\[
\begin{tikzcd}
P^{\bullet} \arrow[r, "\beta"] \arrow[d, "\psi"] & P^{\bullet} \arrow[d, "\psi"] \\
M^{\bullet} \arrow[r, "\alpha"'] & M^{\bullet}
\end{tikzcd}
\]
此外，$R$ 的元素
\[
\operatorname{Tr}(\beta|P^{\bullet}) \overset{\mathrm{def}}{=} \sum_{r} (-1)^{r} \operatorname{Tr}(\beta|P^{r})
\]
不依赖于 $P^{\bullet}$，$\psi$ 与 $\beta$ 的选取.当 $R$-模 $H^{r}(M^{\bullet})$ 都自由时，
\[
\operatorname{Tr}(\beta|P^{\bullet}) = \sum (-1)^{r} \operatorname{Tr}(\alpha|H^{r}(M)).
\]
若 $R$ 是带分式域 $K$ 的整环，则
\[
\operatorname{Tr}(\beta|P^{\bullet}) = \sum (-1)^{r} \operatorname{Tr}(\alpha|H^{r}(M) \otimes_{R} K).
\]
\end{proposition}

\begin{proof}
证明是初等的---见例如 EC VI 13.10.
\end{proof}

下面需要一个复形 $M^{\bullet}$ 存在完美复形 $P^{\bullet}$ 与拟同构 $P^{\bullet} \to M^{\bullet}$ 的判别准则.由于 $H^{r}(P^{\bullet})$ 对所有 $r$ 有限且除有限多个 $r$ 外为零，$H^{r}(M^{\bullet})$ 必须满足同样条件.然而，这些条件不充分：例如，若 $\Lambda = \mathbb{Z}/\ell^{2}\mathbb{Z}$, $M^{0} = \mathbb{Z}/\ell\mathbb{Z}$，其余为 0，则不存在 $P^{\bullet}$，因为不可能在左边截断
\[
\cdots \xrightarrow{\ell} \Lambda \xrightarrow{\ell} \Lambda \xrightarrow{\ell} \cdots \xrightarrow{\ell} \Lambda \to 0 \to \cdots
\]
\[
\parallel
\]
\[
\mathbb{Z}/\ell\mathbb{Z}
\]
在左边截断.

\begin{proposition}\label{prop:29.18}
设 $R$ 是局部 Noether 环，$M^{\bullet}$ 是 $R$-模的复形，使得 $H^{r}(M^{\bullet})$ 对所有 $r$ 有限生成，且对某个 $m$，对 $r > m$ 为零.则存在拟同构 $Q^{\bullet} \to M^{\bullet}$，其中 $Q^{\bullet}$ 是有限生成自由 $R$-模的复形，使得对 $r > m$，$Q^{r} = 0$.若对所有有限生成 $R$-模 $N$，对 $r < 0$，$H^{r}(Q^{\bullet} \otimes_{R} N) = 0$，则存在拟同构 $\psi \colon Q^{\bullet} \to P^{\bullet}$，其中 $P^{\bullet}$ 是 $R$-模的完美复形，使得 $P^{r}$ 只在 $0 \leq r \leq m$ 非零.
\end{proposition}

\begin{proof}
$Q^{\bullet}$ 的存在性是标准结果(见例如 Mumford, \emph{Abelian Varieties}，引理 1, p.47).

考虑序列
\[
Q^{-2} \to Q^{-1} \to Q^{0}.
\]
在对所有 $N$，对 $r < 0$，$H^{r}(Q^{\bullet} \otimes_{R} N) = 0$ 的假设下，它正合，且与 $N$ 作张量积后仍正合.设 $B$ 是 $Q^{-1}$ 在 $Q^{0}$ 中的像.则
\[
Q^{-2} \otimes_{R} N \to Q^{-1} \otimes_{R} N \to B \otimes_{R} N
\]
是复形，且 $Q^{-1} \otimes_{R} N \to B \otimes_{R} N$ 满射(张量积右正合).由此推出 $B \otimes_{R} N \to Q^{0} \otimes_{R} N$ 单射，并把 $B \otimes_{R} N$ 等同于 $Q^{-1} \otimes_{R} N$ 在 $Q^{0} \otimes_{R} N$ 中的像.因此，对 $N$ 的任何子模 $N'$，$B \otimes_{R} N'$ 是 $B \otimes_{R} N$ 的子模.这蕴含 $B$ 平坦.因为 $Q^{0}$ 也平坦，且对每个 $N$，$B \otimes N \to Q^{0} \otimes N$ 单射，推出 $Q^{0}/B$ 平坦(因为对 $r > 0$，$\operatorname{Tor}_r(Q^{0}/B, N) = 0$).定义 $P^{\bullet}$ 为复形
\[
\cdots \to 0 \to Q^{0}/B \to Q^{1} \to Q^{2} \to \cdots.
\]
\end{proof}

注意，若对 $r < 0$，$M^{r} = 0$，则映射 $Q^{\bullet} \to M^{\bullet}$ 诱导映射 $P^{\bullet} \to M^{\bullet}$.

设 $\Lambda = \mathbb{Z}/\ell^{n}\mathbb{Z}$，$E$ 是代数闭域上簇 $U$ 上 $\Lambda$-模的平坦可构造层.设 $j \colon U \hookrightarrow X$ 是 $U$ 的完备化，$j_{!}E \to \mathcal{I}^{\bullet}$ 是 $j_{!}E$ 的内射分解.定理~\ref{thm:15.1},\ref{thm:19.1} 与 K\"unneth 公式表明 $M^{\bullet} \overset{\mathrm{def}}{=} \Gamma(X, \mathcal{I}^{\bullet})$(与 $Q^{\bullet}$)满足命题的假设.我们为任何存在拟同构 $P^{\bullet}(U, E) \to \Gamma(X, \mathcal{I}^{\bullet})$ 的 $\Lambda$-模完美复形写 $P^{\bullet}(U, E)$.于是，$P^{\bullet}(U, E)$ 是表示 $E$ 在 $U$ 上的紧支集上同调的完美复形：
\[
H^{r}(P^{\bullet}(U, E)) = H^{r}_{c}(U, E).
\]
在定理~\ref{thm:29.15} 的情形中，给我们 $H^{r}_{c}(U, E)$ 的自同态 $F$ 的构造也给我们 $\Gamma(X, \mathcal{I}^{\bullet})$ 的自同态 $F$.我们选取 $P^{\bullet}(U, E)$ 上位于 $F$ 之上的自同态 $F$，并定义
\[
\sum (-1)^{r} \operatorname{Tr}(F|H^{r}_{c}(U, E)) = \operatorname{Tr}(F|P^{\bullet}(U, E)).
\]
可以选取随 $n$ 变动的复形 $P^{\bullet}(U, \mathbb{Z}/\ell^{n}\mathbb{Z})$ 形成射影系统，使得映射 $P^{r}(U, \mathbb{Z}/\ell^{n+1}\mathbb{Z}) \to P^{r}(U, \mathbb{Z}/\ell^{n}\mathbb{Z})$ 对所有 $r$ 与 $n$ 诱导同构 $P^{r}(U, \mathbb{Z}/\ell^{n+1}\mathbb{Z})/\ell^{n}P^{r}(U, \mathbb{Z}/\ell^{n+1}\mathbb{Z}) \to P^{r}(U, \mathbb{Z}/\ell^{n}\mathbb{Z})$.我们设 $P^{\bullet}(U, \mathbb{Z}_{\ell})$ 为逆向极限.这是 $\mathbb{Z}_{\ell}$-模的完美复形，使得 $H^{r}(P^{\bullet}(U, \mathbb{Z}_{\ell})) = H^{r}_{c}(U, \mathbb{Z}_{\ell})$.

\subsection*{定理~\ref{thm:29.15} 的证明}

固定有限 Galois 覆盖 $\varphi'_0 \colon V_0 \to U_0$，使得 $E_0|_{V_0}$ 常值，用 $G$ 表示 $V_0$ 在 $U_0$ 上的 Galois 群.则 $E_0$ 对应于带 $G$ 作用的有限自由 $\Lambda$-模 $E$.设 $V = V_0 \otimes_{\operatorname{Specm} \mathbb{F}_q} \mathbb{F}$---它是 $\mathbb{F}$ 上的簇(不一定是连通的)，$\varphi \colon V \to U$ 又是带 Galois 群 $G$ 的 Galois 覆盖.设 $R = \Lambda[G]$.则 $R$ 通过 $G$ 在 $V$ 与 $E$ 上的作用作用于上同调群 $H^{r}(V, E)$.

设 $G_{-1}$ 是满足 $F \circ \varphi = \varphi \circ \alpha$ 的正则映射 $\alpha \colon V \to V$ 的集合.显然，Frobenius 映射 $F \colon V \to V$ 属于 $G_{-1}$，且 $G_{-1} = \{F \circ g \mid g \in G\}$.群 $G$ 通过共轭作用于 $G_{-1}$，对任何 $\alpha \in G_{-1}$，设 $z(\alpha)$ 是 $\alpha$ 在 $G$ 中的中心化子 $Z(\alpha)$ 的阶.

\begin{proposition}\label{prop:29.19}
对所有 $\alpha \in G_{-1}$，$\ell$ 进整数 $\operatorname{Tr}_{\mathbb{Z}_{\ell}}(\alpha|P^{\bullet}(V, \mathbb{Z}_{\ell}))$ 被 $z(\alpha)$ 整除，且
\[
\operatorname{Tr}(F|P^{\bullet}(U, E)) = \sum_{\alpha \in G_{-1}/G} \frac{\operatorname{Tr}_{\mathbb{Z}_{\ell}}(\alpha|P^{\bullet}(V, \mathbb{Z}_{\ell}))}{z(\alpha)} \cdot \operatorname{Tr}_{\Lambda}(F^{-1} \circ \alpha|E)
\]
($\mathbb{Z}/\ell^{n}\mathbb{Z}$ 中的等式).

求和遍历 $G$ 在 $G_{-1}$ 上作用轨道的代表元集 $\alpha$.因为 $F^{-1} \circ \alpha \in G$ 且 $E$ 是有限生成自由 $\Lambda$-模，$\operatorname{Tr}(F^{-1} \circ \alpha|E)$ 是 $\Lambda = \mathbb{Z}/\ell^{n}\mathbb{Z}$ 中的良定义元素.两个迹都不依赖于轨道中 $\alpha$ 的选取.因此右边是 $\Lambda$ 中的良定义元素.
\end{proposition}

现在我们解释为什么该命题完成定理~\ref{thm:29.15} 的证明.如上面所述，可以假定 $U^{F}$ 为空，然后必须证明 $\operatorname{Tr}(F|P^{\bullet}(U, E)) = 0$.为此只需证明对所有 $\alpha$, $\operatorname{Tr}(\alpha|P^{\bullet}(V, \mathbb{Z}_{\ell})) = 0$，但
\[
\operatorname{Tr}(\alpha|P^{\bullet}(V, \Lambda)) \xlongequal{29.17} \sum (-1)^{r} \operatorname{Tr}(\alpha|H^{r}_{c}(V, \mathbb{Q}_{\ell})).
\]
我们希望应用常数系数的 Lefschetz 不动点公式证明 $\sum (-1)^{r} \operatorname{Tr}(\alpha|H^{r}_{c}(V, \mathbb{Q}_{\ell})) = 0$.为此，需要知道
\begin{enumerate}
  \item $\alpha$ 延拓为正则映射 $\alpha \colon Y \to Y$，其中 $Y$ 是含 $V$ 的完备非奇异曲线；
  \item $\alpha$ 在 $Y$ 中的不动点重数为 1；
  \item $\alpha$ 在 $V$ 中没有不动点.
\end{enumerate}
论断 (a) 与 (b) 表明非紧曲线 $V$ 有常数系数的 Lefschetz 不动点公式(比较本节开头的讨论), (c) 表明迹为零.

论断 (a) 是关于曲线及其完备化的一般事实---$\alpha$ 定义 $V$ 上正则函数域 $\mathbb{F}(V)$ 的自同态，它延拓为典范地联系于 $\mathbb{F}(V)$ 的完备非奇异曲线 $Y$ 上的正则映射---见 \S 14.事实上，$G$ 在 $V$ 上的作用唯一延拓为 $Y$ 上的作用，且若 $\alpha = g \circ F$ 在 $V$ 上，则同样的等式在 $Y$ 上成立.对 $Y$ 的任何闭点 $Q$，
\[
(d\alpha)_Q = (dg)_{F(Q)} \circ (dF)_Q,
\]
因为 $(dF)_Q = 0$，它为零.这蕴含 (b)(比较 27.2).

对 (c)，注意 $V$ 中 $\alpha$ 的不动点将位于 $U$ 中 $F$ 的不动点之上，而我们假定 $U^{F} = \varnothing$.

还剩下证明命题~\ref{prop:29.19}.

\begin{proposition}\label{prop:29.20}
(a) 对每个 $N$，复形 $P^{\bullet}(V, \mathbb{Z}/\ell^{N}\mathbb{Z})$ 可以选取为 $(\mathbb{Z}/\ell^{N}\mathbb{Z})[G]$-模的完美复形.对变动的 $N$，它们可以选取形成射影系统，其极限 $P^{\bullet}(V, \mathbb{Z}_{\ell})$ 是 $\mathbb{Z}_{\ell}[G]$-复形，具有对所有 $N$，$P^{\bullet}(V, \mathbb{Z}_{\ell})/\ell^{N}P^{\bullet}(V, \mathbb{Z}_{\ell}) = P^{\bullet}(V, \mathbb{Z}/\ell^{N}\mathbb{Z})$ 的性质.

(b) 有拟同构
\[
P^{\bullet}(U, E) \xrightarrow{\sim} P^{\bullet}(V, \Lambda) \otimes_{R} E.
\]
\end{proposition}

\begin{proof}
(a) 第一个论断的证明是 29.18 之后论证的直接推广.(R 非交换不引起问题；G 相容地作用于 Y 与 E 要求写出这意味着什么的定义.)

(b) K\"unneth 公式(更好的投影公式 22.5)表明
\[
P^{\bullet}(V, E) \xrightarrow{\sim} P^{\bullet}(V, \Lambda) \otimes_{\Lambda} E.
\]
对任何 $G$-模 $P$，迹映射 $\operatorname{Ind}_G(P) \to P$ 诱导同构 $\operatorname{Ind}_G(P)^{G} \to P$.因此
\[
P^{\bullet}(V, E)^{G} \xrightarrow{\sim} P^{\bullet}(U, E),
\]
且
\[
P^{\bullet}(V, E)^{G} \xrightarrow{\sim} (P^{\bullet}(V, \Lambda) \otimes_{\Lambda} E)^{G} = P^{\bullet}(V, \Lambda) \otimes_{R} E.
\]
\end{proof}

\begin{remark}\label{rem:29.21}
设 $N$ 与 $M$ 是 $R$-模.则 $N \otimes_{\Lambda} M$ 是 $R$-模，其中 $\gamma \in G$ 按规则 $\gamma(n \otimes m) = \gamma n \otimes \gamma m$ 作用.设 $M^{0}$ 表示 $G$ 平凡作用的作为 $R$-模的 $M$.则 $r \otimes m \mapsto r \otimes \gamma m \colon R \otimes_{\Lambda} M^{0} \to R \otimes_{\Lambda} M$ 是 $R$-模的同构.因此，若 $M$ 作为 $\Lambda$-模自由，则 $R \otimes_{\Lambda} M$ 作为 $R$-模自由；由此推出若 $N$ 是投射 $R$-模，则 $N \otimes_{\Lambda} M$ 是投射 $R$-模.
\end{remark}

\subsection*{非交换迹插曲}

当 $R$ 是非交换环时，有限秩自由 $R$-模 $M$ 的自同态 $\alpha$ 的迹不是 $R$ 中的良定义元素.例如，假定 $M$ 秩为 1，设 $e$ 与 $e' \overset{\mathrm{def}}{=} be$ 是基元素.若 $\alpha(e) = ae$，则 $\alpha(e') = \alpha(be) = bae$，一般地它不等于 $a(be)$，因为 $ab \neq ba$.设 $R^{\natural}$ 是 $R$ 的加法群除以形如 $ab - ba$ 的元素生成的子群的商.对有限秩自由 $R$-模 $M$ 的自同态 $\alpha$，定义 $\operatorname{Tr}(\alpha|M)$ 为 $\sum a_{ii}$ 在 $R^{\natural}$ 中的像，其中 $(a_{ij})$ 是 $\alpha$ 关于 $M$ 的某组基的矩阵---它不依赖于基的选取.

\begin{lemma}\label{lem:29.22}
设 $M_1$ 与 $M_2$ 是有限秩自由 $R$-模，$\alpha$ 是 $M_1 \oplus M_2$ 的自同态.写 $\alpha = \begin{pmatrix} \alpha_{11} & \alpha_{12}\\ \alpha_{21} & \alpha_{22} \end{pmatrix}$，其中 $\alpha_{ij}$ 是映射 $M_j \to M_i$.则
\[
\operatorname{Tr}_{R}(\alpha|M_1 \oplus M_2) = \operatorname{Tr}_{R}(\alpha_{11}|M_1) + \operatorname{Tr}_{R}(\alpha_{22}|M_2).
\]
\end{lemma}

\begin{proof}
相对于 $M_1$ 的基与 $M_2$ 的基的并计算 $\operatorname{Tr}_{R}(\alpha|M_1 \oplus M_2)$.
\end{proof}

对有限生成投射 $R$-模 $M$ 的自同态 $\alpha$，定义 $\operatorname{Tr}(\alpha|M) = \operatorname{Tr}(\alpha \oplus 0|M \oplus N)$，其中 $N$ 选取为使 $M \oplus N$ 有限秩自由---它不依赖于 $N$ 的选取.

设 $R$ 是群环 $\Lambda[G]$，其中 $\Lambda$ 是交换(Noether)环，$G$ 是有限群.映射 $\sum a_{\gamma}\gamma \mapsto a_e$($e$ 是 $G$ 的单位元)诱导映射 $\varepsilon \colon R^{\natural} \to \Lambda$，对有限生成投射 $R$-模 $M$ 的任何自同态 $\alpha$，定义
\[
\operatorname{Tr}^{\Lambda}_{G}(\alpha|M) = \varepsilon(\operatorname{Tr}_{R}(\alpha|M)).
\]

\begin{lemma}\label{lem:29.23}
对任何有限生成投射 $R$-模 $M$，
\[
\operatorname{Tr}_{\Lambda}(\alpha|M) = [G] \operatorname{Tr}^{\Lambda}_{G}(\alpha|M),
\]
其中 $[G]$ 等于 $G$ 的阶.
\end{lemma}

\begin{proof}
可以假定 $M$ 自由，且由于迹只依赖于 $\alpha$ 矩阵的对角项，假定 $M = R$.则 $\alpha$ 是右乘某个元素 $\sum a_{\gamma}\gamma$ 的作用.对任何 $\delta \in G$, $\alpha(\delta) = a_e\delta + \cdots$，因此 $\operatorname{Tr}_{\Lambda}(\alpha|M) = [G]a_e$.由于 $\operatorname{Tr}^{\Lambda}_{G}(\alpha|M) = a_e$，引理在这种情况下成立.
\end{proof}

这个引理解释了非交换迹对定理~\ref{thm:29.15} 的证明的意义：它们允许把 $\Lambda$ 上的通常迹``除以''群的阶，即使 $\Lambda$ 有限.

\begin{lemma}\label{lem:29.24}
设 $\alpha$ 与 $\beta$ 分别是 $P$ 与 $M$ 的自同态，其中 $P$ 是有限生成投射 $R$-模，$M$ 是作为 $\Lambda$-模自由的有限生成 $R$-模.则
\[
\operatorname{Tr}_{\Lambda}(\alpha \otimes \beta|P \otimes_{\Lambda} M) = \operatorname{Tr}^{\Lambda}_{G}(\alpha|P) \cdot \operatorname{Tr}_{\Lambda}(\beta|M).
\]
\end{lemma}

\begin{proof}
注意，按 (\ref{rem:29.21})，$P \otimes_{\Lambda} M$ 是投射 $R$-模，因此所有项都有定义.只需考虑 $P = R$ 的情形.则 $\alpha$ 是右乘某个元素 $\sum a_{\gamma}\gamma$ 的作用，而 (\ref{rem:29.21}) 的同构 $R \otimes_{\Lambda} M \to R \otimes_{\Lambda} M^{0}$ 把 $\alpha \otimes \beta$ 变为自同态
\[
r \otimes m \mapsto \sum_{\gamma} a_{\gamma} r\gamma \otimes \gamma^{-1}\beta(m)
\]
$R \otimes_{\Lambda} M^{0}$ 的.映射
\[
r \otimes m \mapsto a_{\gamma}r\gamma \otimes \gamma^{-1}\beta(m) \colon R \otimes_{\Lambda} M \to R \otimes_{\Lambda} \check{M}
\]
当 $\gamma = e$ 时迹为 $a_e \operatorname{Tr}_{\Lambda}(\beta|M)$，否则为 0.这完成了证明.
\end{proof}

现在假定有正合序列
\[
1 \to G \to W \xrightarrow{\text{次数}} \mathbb{Z} \to 1.
\]
定义 $G_{-1}$ 为 $W$ 中 $-1$ 的原像.则 $G$ 通过共轭作用于 $G_{-1}$，对 $\alpha \in G_{-1}$，设 $Z(\alpha)$ 是 $\alpha$ 的中心化子，$z(\alpha)$ 是 $\alpha$ 的阶.设 $W^{-}$ 是 $W$ 中映到非正整数 的元素 的集合.设 $P$ 是 $\Lambda$-模，幺半群 $W^{-}$ 在其上 $\Lambda$-线性作用，且作为 $R$-模投射.则 $P^{G}$ 是投射 $\Lambda$-模，且每个 $w \in W^{-}$ 定义 $P^{G}$ 的自同态，只依赖于 $\deg(w)$.

\begin{lemma}\label{lem:29.25}
用上述记号，
\[
\operatorname{Tr}_{\Lambda}\left( \sum_{\alpha \in G_{-1}} \beta|P \right) = \operatorname{Tr}_{\Lambda}\left( \sum_{\alpha \in G_{-1}} \beta|P^{G} \right) = \operatorname{Tr}_{\Lambda}\left( \sum_{\alpha \in G_{-1}} \beta|P_{G} \right).
\]
\end{lemma}

\begin{proof}
固定 $G_{-1}$ 的元素 $\alpha_0$，设 $\eta = \sum_{\gamma \in G} \gamma$.则
\[
P^{G} \hookrightarrow P \xrightarrow{\eta} P^{G}
\]
是乘以 $[G]$.这给出
\[
\operatorname{Tr}_{\Lambda}\left( \sum_{\alpha \in G_{-1}} \alpha|P \right) = \operatorname{Tr}_{\Lambda}(\alpha_0\eta|P) = \operatorname{Tr}([G]\alpha_0|P^{G}) = \operatorname{Tr}\left( \sum_{\alpha \in G_{-1}} \alpha|P^{G} \right)
\]
中的中间等式.乘以 $\eta$ 定义同构 $P_{G} \to P^{G}$(Serre, J.-P., \emph{Corps Locaux}, VIII.1, Prop. 1)，因此 $P^{G}$ 可以换成 $P_{G}$.
\end{proof}

\begin{proposition}\label{prop:29.26}
对任何 $\alpha_0 \in G_{-1}$,
\[
\operatorname{Tr}_{\Lambda}(\alpha_0|P^{G}) = \sum_{\alpha \in G_{-1}/G} \operatorname{Tr}^{\Lambda}_{Z(\alpha)}(\alpha|P).
\]
\end{proposition}

\begin{proof}
见 EC, VI.13.19.
\end{proof}

\subsection*{证明的完成}

还剩下证明命题~\ref{prop:29.19}.
\[
\operatorname{Tr}_{\Lambda}(F|P^{\bullet}(U, E)) = \operatorname{Tr}_{\Lambda}(\varphi^{-1}|P^{\bullet}(V, \Lambda) \otimes_{R} E) \quad \text{(\ref{prop:29.20})}
\]
\[
= \operatorname{Tr}_{\Lambda}(\varphi^{-1}|(P^{\bullet}(V, \Lambda) \otimes_{\Lambda} E)^{G}) \quad \text{(显然)}
\]
\[
= \sum_{\alpha \in G_{-1}/G} \left( \operatorname{Tr}^{\Lambda}_{Z(\alpha)}(\alpha|P^{\bullet}(V, \Lambda) \otimes_{\Lambda} E) \right) \quad \text{(\ref{prop:29.26})}
\]
\[
= \sum_{\alpha \in G_{-1}/G} \operatorname{Tr}^{\Lambda}_{Z(\alpha)}(\alpha|P^{\bullet}(V, \Lambda)) \cdot \operatorname{Tr}_{\Lambda}(F^{-1} \circ \alpha|E) \quad \text{(\ref{lem:29.24})}
\]
\[
= \sum_{\alpha \in G_{-1}/G} \operatorname{Tr}^{\mathbb{Z}_{\ell}}_{Z(\alpha)}(\alpha|P^{\bullet}(V, \mathbb{Z}_{\ell})) \cdot \operatorname{Tr}_{\Lambda}(F^{-1} \circ \alpha|E) \quad \text{(\ref{prop:29.20})}
\]
\[
= \sum_{\alpha \in G_{-1}/G} \frac{\operatorname{Tr}_{\mathbb{Z}_{\ell}}(\alpha|P^{\bullet}(V, \mathbb{Z}_{\ell}))}{z(\alpha)} \cdot \operatorname{Tr}_{\Lambda}(F^{-1} \circ \alpha|E). \quad \text{(\ref{lem:29.23})}
\]
